% !TEX program = xelatex
% Cards format — auto-generated from pioneer-works.tex by cards-convert.mjs
\documentclass[11pt]{article}

\usepackage{fontspec}
\usepackage{unicode-math}
% Latin Modern by OTF filename (not font name): xelatex on macOS (Core
% Text, no fontconfig) cannot resolve "Latin Modern Roman" by name, which
% silently produced ~4KB broken card PDFs. kpathsea finds the OTFs in
% texmf-dist on both macOS and the oven's Linux xelatex.
\setmainfont{lmroman10-regular.otf}[
  BoldFont=lmroman10-bold.otf,
  ItalicFont=lmroman10-italic.otf,
  BoldItalicFont=lmroman10-bolditalic.otf,
]
\setsansfont{lmsans10-regular.otf}[
  BoldFont=lmsans10-bold.otf,
  ItalicFont=lmsans10-oblique.otf,
  BoldItalicFont=lmsans10-boldoblique.otf,
]
\setmonofont{lmmono10-regular.otf}[Scale=0.88,
  BoldFont=lmmonolt10-bold.otf,
  ItalicFont=lmmono10-italic.otf,
]


\usepackage{graphicx}
\graphicspath{{figures/}{../../papers/arxiv-ac/figures/}}
\usepackage{booktabs}
\usepackage{tabularx}
\usepackage{adjustbox}
\usepackage{ragged2e}
\usepackage{microtype}
\usepackage{natbib}



\makeatletter
\def\input@path{{../}}
\makeatother
\usepackage{ac-paper-cards}

% Extra commands from base paper
\newcommand{\papercreated}{2026-05-02}
\newcommand{\paperrevision}{1}
\newcommand{\papersitename}{pioneer-works-dossier-26-arxiv}
\newcommand{\paperurl}{https://papers.aesthetic.computer/\papersitename.pdf}

\hypersetup{
  pdftitle={Pioneer Works: A Dossier},
}

\renewcommand{\acpdfbase}{pioneer-works-dossier-26-arxiv}
\begin{document}

% ============================================================
% TITLE CARD
% ============================================================
\thispagestyle{empty}
\vspace*{0.05em}
\begin{center}
\href{https://papers.aesthetic.computer}{\includegraphics[width=0.72\textwidth]{figures/cover}}\par\vspace{0.25em}
{\acbold\fontsize{18pt}{22pt}\selectfont\color{acdark} Pioneer Works — A Dossier}\par
\vspace{0.1em}
{\fontsize{9pt}{11pt}\selectfont\color{acpink} A Dossier}\par
\vspace{0.2em}
{\scriptsize\color{acgray}\href{https://prompt.ac/@jeffrey}{\color{cyan!70!blue}\textbf{@jeffrey}}\,\textperiodcentered\,Aesthetic{\acdot}Computer\,\textperiodcentered\,ORCID~\href{https://orcid.org/0009-0007-4460-4913}{0009-0007-4460-4913}}\par
\vspace{0.2em}
\rule{0.5\textwidth}{0.5pt}\par
\vspace{0.12em}
\ifacdraft\colorbox{yellow!60}{\small\color{red!80!black}\textbf{\textit{working draft --- not for citation}}}\par
\vspace{0.1em}\fi
{\scriptsize\color{acgray} March 2026 \,\textperiodcentered\, \href{https://github.com/whistlegraph/aesthetic-computer/commit/8fe6526fd}{8fe6526fd}}\par
\end{center}

% ============================================================
% INDEX CARD
% ============================================================
\cardindex

% ============================================================
% BODY
% ============================================================
\section{The Building}

The dossier begins with a building, not an idea. 159 Pioneer Street is a red-brick industrial structure in the Red Hook neighbourhood of Brooklyn, built in 1866 to house the Pioneer Iron Works~\citep{wikipedia2026pioneer} --- a foundry that manufactured steamrollers, railroad track, and machinery for the sugar refineries that lined the Brooklyn waterfront in the late nineteenth century. The street is named after the foundry, not the other way around. The building burned in 1881 and was rebuilt; it operated as light industry through the Second World War and then drifted, like much of Red Hook, into long decline.

In 2011 the artist Dustin Yellin --- a Los Angeles-born sculptor whose ``Frozen Cinema'' glass-collage works had begun selling at six- and seven-figure prices~\citep{wikipedia2026yellin} --- bought the 25{,}000 sq ft shell for \$3.7 million~\citep{bkmag2022pioneer}. He renovated it himself, with his studio crew and the architect Sam Trimble, on an initial budget of roughly \$200{,}000 funded ``largely \dots\ by his own art''~\citep{fvf2017pioneer}. One hundred new windows were cut into the cathedral-scale walls. The space was briefly called \emph{Intercourse}; it opened to the public in 2012. Hurricane Sandy flooded the ground floor with five feet of water that October, and reconstruction had to begin again. The cultural centre re-opened in 2013 under its current name.

\section{The Organization, 2012--2026}

Pioneer Works Art Foundation files as a 501(c)(3) public charity under EIN 46-1097738~\citep{propublica2026pioneer}. The first Form 990 on the federal record (FY2012) reports \$4{,}900 in revenue and \$0 in expenses --- effectively a registration filing for an organization whose physical operations were just beginning.

\subsection{Founding triad (2012--2014)}

Three people anchor the founding period. \textbf{Dustin Yellin} is founder, president, and artist-in-occupancy. \textbf{Gabriel Florenz}, Yellin's second cousin and former studio assistant~\citep{artforum2025stepsdown}, becomes Founding Artistic Director and Vice-President; he ``played a central role in establishing the nonprofit''~\citep{wikipedia2026pioneer} and runs day-to-day artistic programming. In 2014 the cosmologist \textbf{Janna Levin}, a Claire Tow Professor at Barnard / Columbia, joins as Founding Director of Sciences~\citep{wikipedia2026pioneer}. Levin is also a working author (her book on LIGO, \emph{Black Hole Blues}, appears in 2016), and her appointment is the moment the institution becomes legibly multi-disciplinary rather than a generously-staffed artist studio.

\subsection{Programmatic build-out (2014--2019)}

The annual \textbf{Village F\^ete} spring benefit launches 4 May 2014~\citep{pioneerworks2014vf}. Open-call \textbf{residencies} build out across visual arts, music, and technology disciplines. \textbf{Second Sundays} --- a free monthly open house combining open studios, music, and exhibitions --- begins in 2013 and becomes the institution's signature public ritual~\citep{pioneerworks2025ss}. The Amateur Astronomers Association of New York establishes a recurring stargazing presence at Second Sundays under science staff Lauren Silbert~\citep{aaa2018pioneer}. Revenue grows from \$215K (FY2013) to \$14.2M (FY2019, a one-year capital-campaign spike).

\subsection{First professional ED (2019--2022)}

\textbf{Eric Shiner} --- former director of the Andy Warhol Museum (2010--2016), then artistic director at White Cube New York, and earlier SVP for contemporary art at Sotheby's --- is appointed first-ever Executive Director on 19 December 2019, starting 6 January 2020~\citep{prnewswire2019shiner}. His tenure straddles the COVID closure and its aftermath. The \$30 million capital campaign for accessibility renovations and a public observatory is announced in 2019~\citep{archpaper2024reopen}. Shiner departs in 2021; the org runs without an ED through 2022--2023.

\subsection{Mara Manus, capital campaign close (2023--2025)}

\textbf{Mara Manus}, departing executive director of the New York State Council on the Arts (NYSCA, where she had stewarded \$100M in annual funding for seven years), and previously executive director of the Public Theater, becomes inaugural \textbf{Chief Executive Officer} (a new title, distinct from Shiner's ``Executive Director'') in September 2023~\citep{artforum2023manus}. She closes the capital campaign months ahead of schedule and oversees the \$12.7M accessibility renovation~\citep{esd2025renovation}. Pioneer Works re-opens 6 September 2024 after eight months of construction~\citep{archpaper2024reopen}. Manus departs at the end of March 2025 as the org returns to an ``artist- and scientist-led model''~\citep{artforum2025stepsdown}.

\subsection{Florenz era (2025--)}

\textbf{Gabriel Florenz} expands from Founding Artistic Director to \textbf{Executive Director} in March 2025~\citep{artforum2025stepsdown}. \textbf{Stephanie Hemshrot}, who had run capital projects and operations since 2021, is appointed Chief Operating Officer in April 2025~\citep{artnews2025manus}. Levin remains Director of Sciences. Yellin remains President.

\section{People}

\begin{table}[h]
\footnotesize
\centering
\setlength{\tabcolsep}{3pt}
\begin{adjustbox}{max width=\linewidth,center}%
\begin{tabularx}{\linewidth}{lXX}
\toprule
\textbf{Tenure} & \textbf{Top exec.} & \textbf{Title} \\
\midrule
2012-- & Dustin Yellin & President \& Founder \\
2012-- & Gabriel Florenz & Artistic Dir. \\
2014-- & Janna Levin & Dir. of Sciences \\
2020--2021 & Eric Shiner & Executive Dir. \\
2023--2025 & Mara Manus & CEO \\
Mar 2025-- & Gabriel Florenz & Executive Dir. \\
Apr 2025-- & Stephanie Hemshrot & COO \\
\bottomrule
\end{tabularx}%
\end{adjustbox}
\caption{Top-of-house leadership of record. Yellin and Florenz are second cousins; Florenz was Yellin's studio assistant before becoming Founding Artistic Director.}
\label{tab:eds}
\end{table}

The board is co-led, in the FY2024 Form 990 era, by chair \textbf{Austin Hearst} (Hearst family; named chair in 2023~\citep{artforum2023manus}) and co-vice-chair \textbf{David Belt}. Other recurring or recently named directors: \textbf{Julia Bator}, \textbf{Alan Tisch}, \textbf{Andres Santo Domingo}, \textbf{Leyli Zohrenejad}, \textbf{Kyra Tirana Barry}, \textbf{Alex Black}, \textbf{Mary Cunney}, \textbf{Mark Dalio}, \textbf{Tiffiney Davis}, \textbf{Ron Elum}, \textbf{Matthew Harris}, \textbf{Stephanie Ingrassia}, \textbf{Neal Katyal}, \textbf{Tom X.~Lee}, \textbf{Charles Meyer}, \textbf{Abby Pucker}, and \textbf{Richard Sarnoff}~\citep{prnewswire2019shiner,artforum2023manus}. The 2022 anniversary press cycle named an advisory cohort that included \textbf{David Byrne}, \textbf{Princess Eugenie of York}, \textbf{Justin Vernon} (Bon Iver), and \textbf{Kehinde Wiley}~\citep{bkmag2022pioneer}.

The science department head from the early years was the nano-physicist \textbf{Matthew Putman}, hosting the long-running \emph{Scientific Controversies} discussion series alongside Levin~\citep{wikipedia2026pioneer}. \textbf{Joshua Jelly-Schapiro} co-edits \emph{Broadcast} with Levin; \textbf{Lucy McKeon} and \textbf{David Everitt Howe} are senior editors; \textbf{Michael Jones} is co-founding managing editor~\citep{starrevue2024broadcast}.

\section{Programs}

\textbf{Residencies} (2013--) --- The institution's central commitment. Three open-call tracks have run for most of the period:

\begin{itemize}
\item \emph{Visual Arts}, 5 months, \$2{,}500 honorarium, 12 residents/year, public-facing glass-walled studios~\citep{artdeadline2026residency}.
\item \emph{Music}, 1 month, \$1{,}000 honorarium, 12 residents/year, exclusive use of a professional recording studio~\citep{artdeadline2026residency}.
\item \emph{Technology}, $\sim$3 months, formerly open-call and now invite-only; tied to the on-site Tech Lab (3D printers, laser cutter, VR, robotics, AI, additive manufacturing)~\citep{pioneerworks2025tech}.
\end{itemize}
A \emph{Narrative Arts} track was piloted in the 2020 open call~\citep{hyperallergic2019openall}. Residents have included Baseera Khan, Jen Liu, Kameelah Janan Rasheed, Martine Syms, Morehshin Allahyari, Mary Helena Clark, Caroline Sinders, Xin Liu, Avram Finkelstein, Jess Johnson, LJ Roberts, Alexandria Smith, Erin Johnson, Nate Lewis, MSHR, and Miho Hatori, among many others (year-by-year list in \texttt{data/residents.csv}).

\textbf{Science} (2014--) --- Discussion series \emph{Scientific Controversies} (founded with Putman + Levin), one-off lectures, and from 2024 on the planned \emph{Public Observatory} on the building's roof: a 16-foot restored 1890s antique telescope to be operated as ``NYC's first free public observatory''~\citep{surface2024obs}. The science programme is the through-line of the Levin years and is also what reliably draws non-art press.

\textbf{\emph{Broadcast}} (2020--) --- Online magazine of arts and sciences, edited by Levin and Joshua Jelly-Schapiro. Print Issue 01 in November 2023, featuring hannah baer, Janna Levin, Ed Park, and Rowan Ricardo Phillips~\citep{starrevue2024broadcast}. Replaces the earlier in-house publication \emph{Intercourse} (the same word Yellin briefly used for the building) and the pamphlet series \emph{Groundworks} and \emph{MIMO}~\citep{fvf2017pioneer}.

\textbf{Pioneer Works Press} (\textasciitilde 2016--) --- Imprint for ``experimental and pathbreaking work from leading artists and writers in contemporary culture''; co-distributed with Printed Matter and others.

\textbf{Second Sundays} (2013--) --- Free monthly open house, second Sunday of every month. Open studios with current residents, music, exhibitions, food, garden access. Hosts hundreds per event~\citep{aaa2018pioneer}.

\textbf{Open House / OPEN} (2024 reopening + ongoing) --- The expanded post-renovation public-facing programme; the October 2024 \emph{Fall Benefit} re-opening was branded ``OPEN''~\citep{surface2024gala}.

\textbf{Village F\^ete} (May 2014--) --- Annual spring benefit gala. Inaugural 4 May 2014; the 2025 edition (6 May 2025) featured David Byrne~\citep{artnews2025fete}. The 7th annual was cancelled (COVID, 2020). The fete is the institution's principal annual fundraiser; Pioneer Works also runs an associated benefit auction (\texttt{auction.pioneerworks.org}).

\textbf{Education} (ongoing) --- Free workshops for K-12 students through partnerships with the Red Hook Initiative and others. 300+ K-12 students reached in 2022 alone~\citep{bkmag2022pioneer}.

\section{Money}\label{sec:money}

The federal tax record (Pioneer Works Art Foundation, EIN 46-1097738~\citep{propublica2026pioneer}) is the cleanest source. Filings are public on ProPublica's Nonprofit Explorer back to FY2012, the foundation's first year.

\subsection{Top-line, FY2012--FY2024}

\begin{table*}[t]
\footnotesize
\centering
\setlength{\tabcolsep}{4pt}
\begin{adjustbox}{max width=\linewidth,center}%
\begin{tabularx}{\linewidth}{Xrrrrrr}
\toprule
\textbf{FY end} & \textbf{Revenue} & \textbf{Contri.\,/\,Grants} & \textbf{Prog.\,svc.} & \textbf{Expenses} & \textbf{Top exec. comp.} & \textbf{Net assets} \\
\midrule
2012 & 4{,}900 & 4{,}900 & --- & 0 & --- & 4{,}900 \\
2013 & 215{,}291 & 182{,}500 & --- & 168{,}404 & --- & 51{,}787 \\
2014 & 1{,}306{,}327 & 886{,}624 & 362{,}356 & 1{,}171{,}643 & --- & 186{,}471 \\
2015 & 2{,}869{,}041 & 1{,}989{,}606 & 500{,}318 & 2{,}283{,}062 & --- & 772{,}450 \\
2016 & 3{,}625{,}522 & 2{,}640{,}303 & 758{,}878 & 3{,}469{,}242 & --- & 928{,}730 \\
2017 & 6{,}055{,}241 & 4{,}941{,}234 & 910{,}632 & 5{,}544{,}113 & 150{,}000 & 1{,}439{,}858 \\
2018 & 7{,}506{,}234 & 6{,}723{,}647 & 721{,}731 & 6{,}680{,}659 & 162{,}500 & 2{,}265{,}433 \\
2019 & \textbf{14{,}225{,}980} & 12{,}184{,}712 & 1{,}644{,}602 & 7{,}557{,}149 & 144{,}406 & 8{,}934{,}264 \\
2020 & 9{,}737{,}121 & 9{,}280{,}986 & 453{,}450 & 5{,}817{,}051 & 286{,}731 & 12{,}854{,}334 \\
2021 & 7{,}041{,}619 & 6{,}659{,}196 & 381{,}460 & 6{,}535{,}212 & 233{,}140 & 13{,}360{,}741 \\
2022 & 10{,}424{,}515 & 9{,}590{,}229 & 833{,}744 & 8{,}666{,}931 & 188{,}928 & 15{,}118{,}325 \\
2023 & \textbf{13{,}401{,}014} & 11{,}709{,}988 & 970{,}411 & 10{,}555{,}334 & 216{,}000 & 17{,}864{,}005 \\
2024 & 9{,}963{,}860 & 9{,}059{,}301 & 606{,}712 & 8{,}229{,}322 & 299{,}738 & 19{,}598{,}543 \\
\bottomrule
\end{tabularx}%
\end{adjustbox}
\caption{Pioneer Works Art Foundation --- top-line Form 990 figures by fiscal year (USD; FY end is calendar year-end based on filing pattern). Source: ProPublica Nonprofit Explorer for EIN 46-1097738~\citep{propublica2026pioneer}. Government-grants line is not separately tabulated in the ProPublica summary --- detail TBD from the IRS XML pull. Top-exec. comp. is the highest single officer line for each year (titles vary across the period: Florenz, Shiner, Manus depending on year).}
\label{tab:financials}
\end{table*}

The shape, in plain English: revenue is contribution-driven (88--96\% in nearly every year), with program-service revenue (event rentals, workshop fees, ticket sales) running \$0.4--\$1.6M but never the dominant line. Two years break the trend upward: \textbf{FY2019} at \$14.2M revenue (the year of the \$30M capital-campaign launch and a single-year contributions surge to \$12.2M) and \textbf{FY2023} at \$13.4M (closing-year of the campaign, with the \$12.7M renovation work funded). Net assets grew almost monotonically from \$5K (FY2012) to \$19.6M (FY2024). The institution by FY2024 is, on the federal record, an order of magnitude larger than Rhizome (compare \$821K revenue, \$262K net assets for Rhizome FY2024).

\subsection{The Yellin-as-funder mechanism}

The founder-as-funder dynamic is unusual enough to warrant a paragraph as a fact about the funding model. The 2011 building purchase (\$3.7M) and the initial 2011--2012 renovation budget (``largely \dots\ funded by his own art''~\citep{fvf2017pioneer}) came directly from Yellin's art sales. Press coverage of Pioneer Works' first decade consistently treats this as a known: ``initially financed through Yellin's artwork sales''~\citep{fvf2017pioneer}; ``initial budget of \$200{,}000 that was largely funded by his own art''~\citep{fvf2017pioneer}. Once the foundation incorporated as a 501(c)(3) in 2012, the institution moved progressively onto a mixed funding model: institutional grants, foundation gifts, ticketed programming, the annual Village F\^ete, individual donors. Yellin remains President, and the building remains the foundation's principal physical asset; the property-ownership relationship between Yellin (as private acquirer) and the foundation (as occupant / owner of record) over time is not fully recoverable from the public 990s and is flagged in \texttt{data/README.md}.

\subsection{Capital campaign and accessibility renovation}

The 2019 capital campaign was set at \textbf{\$30 million}~\citep{archpaper2024reopen}. By August 2024 \$28.3M had been raised; the remaining \$1.7M was earmarked for the public observatory. The renovation itself was \$12.7M of construction work~\citep{esd2025renovation}, of which \$2.3M came from New York State (Empire State Development) public funds; the balance \$10.4M was private. The building was closed for eight months in 2024 and re-opened on 6 September 2024.

\subsection{Annual gala}

The \textbf{Village F\^ete} (May, annual since 2014; 7th annual cancelled 2020 for COVID) is the principal annual fundraiser; Pioneer Works also runs an associated online \emph{Benefit Auction}. A separate \textbf{Fall Benefit / OPEN} re-opening gala occurred 8 October 2024~\citep{surface2024gala}. The 2025 Village F\^ete brought David Byrne to Red Hook on 6 May 2025~\citep{artnews2025fete}. Per-event Schedule G (gross, contributions, direct expenses, net) figures are deferred --- the IRS XMLs were not pulled for this first pass. Schedule G detail is recorded in \texttt{data/README.md} as a TBD.

\section{Footprint}

\textbf{The building.} 159 Pioneer Street, Red Hook, Brooklyn. 25{,}000 sq ft of indoor space (some sources say 24{,}000 or 27{,}000 sq ft; the post-2024 renovation puts the figure at 25{,}000 plus a 20{,}000 sq ft adjoining garden, with two new mezzanines and an ADA elevator). 40-foot ceilings. 100 windows from the 2011 renovation. Originally built 1866 for Pioneer Iron Works; the street is named after the foundry. Owned and operated by the foundation.

\textbf{The neighbourhood.} Red Hook in 2012 was an under-served, underbuilt waterfront industrial neighbourhood with poor subway access (the F train at Smith--9th Street is the nearest). Pioneer Works' presence is large enough to function as a neighbourhood anchor: Second Sundays draw hundreds, and Pioneer Works partners with the Red Hook Initiative on K-12 programming. The Hurricane Sandy flood (October 2012) was a neighbourhood event before it was an institutional one.

\textbf{The aesthetic register.} Pioneer Works' identifying aesthetic, in press coverage of all eras, is the \emph{cathedral-of-process} register: cast-iron columns, exposed brick, glass-walled artist studios, scientists at chalkboards beside artists at workstations, Yellin's own glass-collage sculptures often in the background. The institution names itself a ``social sculpture'' and a ``museum of process''~\citep{fvf2017pioneer}.

\textbf{The cultural footprint.} Recurring partners include Columbia / Barnard (via Levin), the Amateur Astronomers Association of New York, the Red Hook Initiative, Printed Matter, Empire State Development (post-2025 renovation press), and a long advisory list including David Byrne, Justin Vernon, and Kehinde Wiley. \emph{Broadcast} (2020--) is the institution's contribution to the periodicals shelf. Pioneer Works Press is its book imprint.

\section{Reading list}

\subsection{Press / coverage --- already cited}

\textbf{ArtNews, ``After Ten Years, Pioneer Works\dots''} (2022)~\citep{artnews2022ten} --- the anniversary recap.\\
\textbf{BkMag, ``Pioneers at Work''} (2022)~\citep{bkmag2022pioneer} --- decade-end community piece.\\
\textbf{Brooklyn Eagle / DNyuz / ArtForum on Manus appointment} (Sept 2023)~\citep{artforum2023manus}.\\
\textbf{ArtForum on Manus departure} (May 2025)~\citep{artforum2025stepsdown}.\\
\textbf{Surface Mag} on the renovation~\citep{surface2024obs} and on the OPEN gala~\citep{surface2024gala}.\\
\textbf{Architect's Newspaper} on ADA reopening~\citep{archpaper2024reopen}.\\
\textbf{Empire State Development press release} on the \$12.7M renovation~\citep{esd2025renovation}.\\
\textbf{Friends of Friends / FvF} on the founding-decade aesthetic~\citep{fvf2017pioneer}.

\subsection{To acquire / read deeper}

\textbf{Joseph Beuys, \emph{What Is Art?}} --- on social sculpture, the term Pioneer Works applies to itself.\\
\textbf{Janna Levin, \emph{Black Hole Blues and Other Songs from Outer Space}} (2016) --- the Director of Sciences' principal book of the institutional period.\\
\textbf{Dustin Yellin oral history / interviews} --- including the Friends of Friends profile and the various BkMag and ArtNews recaps.\\
\textbf{The Red Hook neighbourhood literature} --- on post-Sandy infrastructure, gentrification arguments, and the relationship between cultural-anchor institutions and waterfront industrial neighbourhoods.\\
\textbf{\emph{Broadcast} Issue 01} (Nov 2023) --- the institution's own most extended self-presentation.\\
\textbf{Comparative residency programme literature} --- Skowhegan, MacDowell, Yaddo, Headlands, Fountainhead, NEW INC; Pioneer Works' open-call, multi-disciplinary, urban-anchor format is unusual in the residency landscape.

\section{What this dossier is not}

\begin{itemize}
  \item Not a thesis. There is no question being answered.
  \item Not a critique. The Yellin-as-founder-and-funder mechanic is recorded as a fact about the institution's funding model, neither praised nor condemned.
  \item Not a comparison. Adjacent organizations (Rhizome, NEW INC, Eyebeam, Knockdown Center, MASS MoCA, the Watermill Center) are mentioned only where they directly intersect Pioneer Works or where the contrast in scale clarifies the federal-tax-record numbers.
  \item Not exhaustive. The dossier covers what is publicly recoverable; gaps (Schedule G auction detail, full year-by-year residents, named-grant inventory below the seven-figure level) are flagged inline. The companion data folder (\texttt{papers/arxiv-pioneer-works/data/}) contains the underlying CSVs.
\end{itemize}

\vspace{0.5em}
\noindent\textbf{ORCID:} \href{https://orcid.org/0009-0007-4460-4913}{0009-0007-4460-4913}

\bibliographystyle{plainnat}
\bibliography{references}

\end{document}
